\documentclass[11pt]{article}
\usepackage[margin=1in]{geometry}
\usepackage{booktabs,longtable,array,graphicx,enumitem}
\usepackage[hidelinks]{hyperref}
\setlength{\parindent}{0pt}
\setlength{\parskip}{6pt}
\renewcommand{\arraystretch}{1.15}
\newcommand{\plainbox}[1]{\par\vspace{2pt}\noindent\fbox{\parbox{\dimexpr\linewidth-2\fboxsep-2\fboxrule\relax}{#1}}\par\vspace{2pt}}
\begin{document}
{\Large\textbf{Tasks Module, Part 2: A Test Dataset}}\par
{\large How the made-up answers were built, and what they can and cannot tell us}\par\vspace{4pt}\hrule\vspace{8pt}

\section*{1. Why we made a test dataset}
We have not run the survey yet. Before we do, we want to know that the answers to the task questions
can be turned into useful results. So we made up answers for 200 imaginary workers and ran the
analysis on them. This is the same idea as the earlier test data for the poverty work.

\plainbox{\textbf{Important.} Every pattern in this data comes from a rule we wrote down. Nothing in it is a finding
about real workers. It tests the \emph{questions and the method}, not the world.}

\section*{2. What we started from}
We reused the imaginary workers already made for the poverty work: 200 people, filled by fixed quotas for each job.
After the same drop-out as before, 104 remain (the ones that would be ``interviewed''). We kept their
job, off-season work, age, schooling, place of origin, phone and training. The task answers were then made to fit those facts.

{\small
\begin{longtable}{lrr}
\toprule
\textbf{Job} & \textbf{Full pool} & \textbf{After drop-out} \\
\midrule
\endfirsthead
\toprule
\textbf{Job} & \textbf{Full pool} & \textbf{After drop-out} \\
\midrule
\endhead
Pony worker & 22 & 5 \\
Shop owner & 20 & 11 \\
Shopkeeper & 20 & 14 \\
Wage labourer & 20 & 11 \\
Hotel/lodging worker & 18 & 15 \\
Porter & 18 & 10 \\
Hotel/lodging owner & 18 & 7 \\
Dhaba/food-stall owner & 12 & 6 \\
Driver & 12 & 7 \\
Other & 12 & 4 \\
Pony business & 12 & 5 \\
Palki & 12 & 6 \\
Guide & 4 & 3 \\
\textbf{All} & \textbf{200} & \textbf{104} \\
\bottomrule
\end{longtable}
}

\section*{3. How the answers were made}
\textbf{Step 1. Who does which task, by job.} For each of the 13 jobs we sorted the 27 tasks into four levels:
\emph{core} tasks (a worker does it regularly with chance 90 in 100), \emph{common} tasks (55 in 100),
\emph{occasional} tasks (20 in 100) and \emph{rare} tasks (3 in 100).
The core tasks come from the wording of the job in India's National Classification of Occupations (NCO 2015).
Where NCO has no entry (palki and dandi bearers) they are our judgement, and the last column of the table in section 4 says so.

\textbf{Step 2. People differ.} Some people simply do more of everything. More years in the work adds a little.
More schooling adds a little to tasks that use records and judgement.

\textbf{Step 3. Work done before elsewhere (answer 2).} The chance starts low (6 in 100).
It rises with off-season work: farming raises animal care and carrying; wage labour raises building and loading;
petty trade raises selling and cash. A past job in a different trade brings in that trade's tasks. Older workers have a little more.

\textbf{Step 4. Main task.} It is the job's usual main task, if the person does it regularly.
Otherwise it is the person's most typical regular task.

\textbf{Step 5. Job conditions and skills.} Made from the person's job, schooling, place of origin, phone and training.
For example, someone who supervises others regularly must report at least one person directed.

\section*{4. The core tasks we gave each job}
{\footnotesize
\begin{longtable}{p{2.7cm}p{5.6cm}p{6.5cm}}
\toprule
\textbf{Job} & \textbf{Core tasks} & \textbf{Where it comes from} \\
\midrule
\endfirsthead
\toprule
\textbf{Job} & \textbf{Core tasks} & \textbf{Where it comes from} \\
\midrule
\endhead
Pony worker & 5 lead or control pack animals; 3 load, unload, stack, count goods; 6 feed or care for animals & NCO 9332.0400 Pack Animal Driver: loads and unloads animals, leads them, feeds them, cleans stables. \\
Pony business & 5 lead or control pack animals; 14 sell goods or services; 15 bargain over prices or fares; 16 handle cash and payments; 26 supervise or hire others & NCO 9332.0400 plus owner tasks (selling rides, cash, hiring handlers) from NCO 1120.2000 retail proprietor logic. \\
Porter & 1 carry loads on foot; 3 load, unload, stack, count goods & NCO 9621.9900 luggage porters and 9333.0100 loader and unloader. \\
Palki & 2 carry a passenger; 1 carry loads on foot & *No NCO code. Judgement: carries passengers and loads; needs union check. \\
Shopkeeper & 14 sell goods or services; 4 pack or tie goods; 20 clean rooms or public areas & NCO 5223.0100 Shop Assistant: measures and packs goods, delivers, keeps shop tidy, sells. \\
Shop owner & 14 sell goods or services; 15 bargain over prices or fares; 16 handle cash and payments; 17 buy stock or arrange supplies; 25 keep records or accounts & NCO 1120.2000 Working Proprietor, Retail Trade: buys merchandise and sells it for profit. \\
Hotel/lodging worker & 19 serve guests; 20 clean rooms or public areas & NCO 5131.0200 Steward Hotel, 5151.0800 Room Bearer, 9112.9900 hotel cleaners. \\
Hotel/lodging owner & 14 sell goods or services; 16 handle cash and payments; 17 buy stock or arrange supplies; 25 keep records or accounts; 26 supervise or hire others & NCO 5151.0600 Hotel and Restaurant Keeper: buys supplies, appoints staff, serves guests, collects money. \\
Wage labourer & 1 carry loads on foot; 3 load, unload, stack, count goods; 11 build or construct & NCO 9313.9900 Building Construction Labourers: simple routine tasks in construction. \\
Other & none (mixed group) & Mixed group; no single NCO code. \\
Driver & 7 drive a motor vehicle; 16 handle cash and payments & NCO 8322.0501 Light Motor Vehicle Driver: drives safely on an assigned route; fares from NCO cashier logic. \\
Dhaba/food-stall owner & 14 sell goods or services; 16 handle cash and payments; 17 buy stock or arrange supplies; 18 cook or prepare food or drink & NCO 5212.9900 Street Food Vendors (Dhabawala) and 5151.0600 Hotel and Restaurant Keeper. \\
Guide & 21 guide or explain to visitors & NCO 5113.0200 Tourist Guide: guides visitors, explains sites, answers questions, may assist shopping. \\
\bottomrule
\end{longtable}
}

\section*{5. Checks on the made-up data}
We ran 11 checks. All 11 passed. The checks catch answers that could not happen in real life.
\begin{itemize}[leftmargin=1.2em,itemsep=1pt]
\item All 27 task answers are 1, 2 or 3 with no missing.
\item Every worker has at least one regular task.
\item Main task is a task coded 1.
\item People directed (t2\_07) is at least 1 exactly when 'supervise or hire' (t1\_26) is coded 1.
\item Trips per day is asked only when a carrying or animal task is coded 1.
\item Course type recorded exactly when a course was taken.
\item Digital payments imply a smartphone and internet.
\item Read/write in a language implies speaking it.
\item Writing at work implies more than chance of reading (no writer without reader more than 12\%).
\item Palki bearers do the 'carry a passenger' task regularly (most of them).
\item Drivers hold a car licence (most of them).
\end{itemize}
The same checks can be run on the real survey data as a first quality test.

\section*{6. What the test data can and cannot show}
\textbf{It can show} whether the answers fit together, whether the analysis code runs, and what kind of result each method gives.
\textbf{It cannot show} what Kedarnath workers really do. Every job profile is our assumption, and the sample is small
(104 people). Two results that look different in this data may not differ in real life.
When real answers arrive, the code is run again on them without change.

\section*{7. Files}
\begin{itemize}[leftmargin=1.2em,itemsep=1pt]
\item \texttt{02\_data/generate\_task\_data.py}: the code that makes the data.
\item \texttt{02\_data/tasks\_synth\_full.csv} and \texttt{tasks\_synth\_fielded.csv} (also \texttt{.dta} for Stata).
\item \texttt{02\_data/occupation\_profiles.csv}: the levels given to each job.
\end{itemize}
\end{document}
